
\medskip
\noindent
\SGARefTarget{sga7.sga7i.eix.n7-2-2}{7.2.2}. We apply \SGARefLink{sga7.sga7i.eix.n7-2-1}{7.2.1} to the formal group scheme $\widehat A$
considered in \SGARefLink{sga7.sga7i.eix.n7-1}{7.1}, in the favorable case treated there in which the
formal abelian scheme $\widehat B$ is algebraizable (cf. \SGARefLink{sga7.sga7i.eix.n7-1-5}{7.1.5}), and
we suppose in addition that the torus $T_o$ over $S_o$ is isotrivial.
We thereby obtain a group scheme $A^\natural$ over $S$, called the
\underline{Raynaud group} of $A$ over $S$, characterized up to
isomorphism by a canonical isomorphism of formal group schemes over
$\operatorname{Spf}(V)=\widehat S$:
\begin{equation}\tag{\SGARefTarget{sga7.sga7i.eix.n7-2-3}{7.2.3}}
 (A^\natural)^\wedge\simeq\widehat A.
\end{equation}
It follows from \SGARefLink{sga7.sga7i.eix.n7-2-1}{7.2.1} that
\begin{equation}\tag{\SGARefTarget{sga7.sga7i.eix.n7-2-4}{7.2.4}}
 \Phi=A^\natural/A^{\natural o}
\end{equation}
is a finite etale group over $S$, and that the identity component
\[
 A^{\natural o}=A^{o\natural}
\]
is endowed with the structure of an extension
\begin{equation}\tag{\SGARefTarget{sga7.sga7i.eix.n7-2-5}{7.2.5}}
 0\longrightarrow T\longrightarrow A^{\natural o}
 \longrightarrow B\longrightarrow0,
\end{equation}
where $B$ is the abelian scheme over $S$ whose formal completion is
the formal abelian scheme $\widehat B$ constructed in \SGARefLink{sga7.sga7i.eix.n7-1}{7.1}, and $T$ is
characterized, up to a unique isomorphism, as the isotrivial torus
over $S$ lifting $T_o$.  On restriction to the $S_n$, the extension
structure (\SGARefLink{sga7.sga7i.eix.n7-2-5}{7.2.5}) gives the extension structures on
$A_n^{\natural o}=A_n^o$ defined by (\SGARefLink{sga7.sga7i.eix.n7-1-1}{7.1.1}).

\medskip
\noindent
\SGARefTarget{sga7.sga7i.eix.n7-2-6}{7.2.6}. It also follows immediately from \SGARefLink{sga7.sga7i.eix.n7-2-1}{7.2.1} that the Raynaud group
$A^\natural$, as well as the exact sequence (\SGARefLink{sga7.sga7i.eix.n7-2-5}{7.2.5}), depends
functorially on $A$ (satisfying the conditions under consideration),
and that its formation commutes with every base change
$S'\longrightarrow S$ associated with a homomorphism of adic
topological rings $V\longrightarrow V'$.
